% Compile with pdflatex, bibtex, pdflatex, pdflatex.
\documentclass[11pt,a4paper]{article}
\usepackage[T1]{fontenc}
\usepackage[utf8]{inputenc}
\usepackage{lmodern,microtype}
\usepackage[margin=27mm]{geometry}
\usepackage{amsmath,amssymb,amsthm,mathtools,mathrsfs,booktabs}
\usepackage[hidelinks]{hyperref}
\usepackage{bookmark}
\newtheorem{theorem}{Theorem}[section]
\newtheorem{lemma}[theorem]{Lemma}
\newtheorem{proposition}[theorem]{Proposition}
\newtheorem{corollary}[theorem]{Corollary}
\newtheorem{conjecture}[theorem]{Conjecture}
\theoremstyle{remark}
\newtheorem{remark}[theorem]{Remark}
\numberwithin{equation}{section}
\DeclareMathOperator{\Hom}{Hom}
\DeclareMathOperator{\Tr}{Tr}
\DeclareMathOperator{\diag}{diag}
\newcommand{\E}{\mathbb E}
\newcommand{\Prb}{\mathbb P}
\newcommand{\R}{\mathbb R}
\newcommand{\Q}{\mathbb Q}
\newcommand{\dd}{\,\mathrm d}
\newcommand{\one}{\mathbf 1}
\newcommand{\ip}[2]{\langle #1,#2\rangle}
\newcommand{\op}{\mathrm{op}}
\newcommand{\HS}{\mathrm{HS}}
\newcommand{\maj}{\mathcal P}
\newcommand{\code}[1]{\texttt{\detokenize{#1}}}
\setlength{\emergencystretch}{3em}
\hypersetup{pdftitle={Commonness of independent apices of trees and even cycles},
 pdfauthor={Taeyoung Kim},
 pdfsubject={Ramsey multiplicity, entropy, and exact flag algebra certificates},
 pdfkeywords={common graphs, independent apices, trees, even cycles, entropy, flag algebras}}
\title{Commonness of independent apices of trees and even cycles}
\author{Taeyoung Kim}
\date{October 2, 2026}
\begin{document}
\maketitle

\begin{abstract}
Let $H^{+k}$ denote the graph obtained from $H$ by adjoining an
independent set of $k$ vertices, each joined to every vertex of $H$.
We prove that $H^{+k}$ is common for every positive integer $k$ when
$H$ is a tree or an even cycle. These results establish two infinite
families of the independent-apex conjecture of Grzesik, Lee,
Lidick\'y, and Volec. For tree bases, we obtain a sharp density
inequality by extending the law of a random book along the tree and
estimating its relative entropy against the homomorphism weight. For even
cycles, four exact six-vertex flag algebra certificates yield a
weighted fourth-moment comparison. Spectral inequalities and
convexity then extend this comparison to all even cycle lengths and
all apex numbers. For at least three apices, we obtain a stronger
comparison with the monochromatic density of the base cycle. Equality
in the commonness bound for an apex of an even cycle holds exactly
at the constant graphon $1/2$.
\end{abstract}

\smallskip
\noindent\textit{2020 Mathematics Subject Classification.}
Primary 05C35; secondary 05C60, 94A17.
\par\smallskip
\noindent\textit{Keywords.} Ramsey multiplicity, common graphs,
independent apices, entropy, flag algebras.

\section{Introduction}\label{sec:intro}

The Ramsey multiplicity problem asks how few monochromatic copies
of a fixed graph can occur in a two-coloring of a large complete
graph. An independent fair coloring supplies a natural benchmark:
a graph with $e\geq1$ edges is monochromatic with probability $2^{1-e}$.
A graph is called \emph{common} if no coloring improves this
benchmark asymptotically. Goodman's theorem shows that the triangle
is common \cite{Goodman}, whereas Thomason proved that every
complete graph of order at least four is uncommon \cite{Thomason}.
The resulting distinction makes it natural to seek graph operations
that preserve commonness.

We work with graphons. For a symmetric measurable function
$W:[0,1]^2\to[0,1]$ and a finite simple graph $F$, define
\[
 t(F,W)=\int_{[0,1]^{V(F)}}\prod_{uv\in E(F)}W(x_u,x_v)
                    \prod_{u\in V(F)}\dd x_u,
 \qquad M(F,W)=t(F,W)+t(F,1-W).
\]
If $e(F)=|E(F)|$, commonness is equivalent to
\begin{equation}\label{eq:common}
 M(F,W)\geq2^{1-e(F)}\qquad\text{for every graphon }W.
\end{equation}
The constant graphon $W\equiv1/2$ attains equality. A direct
comparison with finite colorings is given in
Appendix~\ref{app:finite}.

For $k\geq1$, write $H^{+k}$ for the graph with vertex set
$V(H)\sqcup\{z_1,\ldots,z_k\}$ and edge set
\[
 E(H)\cup\{vz_i:v\in V(H),\ 1\leq i\leq k\}.
\]
The added vertices are pairwise nonadjacent; we call them independent
apices. This operation connects commonness with Sidorenko's
conjecture. A bipartite graph $H$ has the \emph{Sidorenko property}
if
\[
 t(H,W)\geq t(K_2,W)^{e(H)}\qquad\text{for every graphon }W.
\]
Sidorenko's conjecture asserts this property for every bipartite
graph \cite{SidorenkoCorrelation}. Applying it to both colors and
using convexity immediately gives commonness of $H$. Sidorenko
proved the stronger fact that $H^{+1}$ is common whenever $H$ is
connected and has the Sidorenko property \cite{SidorenkoFriendly}.

Grzesik, Lee, Lidick\'y, and Volec proposed extending this result to
arbitrarily many independent apices \cite[Conjecture~1.3]{GLLV}.

\begin{conjecture}[Grzesik--Lee--Lidick\'y--Volec]\label{conj:apices}
If $H$ is a connected bipartite graph with the Sidorenko property,
then $H^{+k}$ is common for every positive integer $k$.
\end{conjecture}

Their paper verifies the conjecture for connected bipartite bases
on at most five vertices \cite[Theorem~1.4]{GLLV}. It also treats
the graphs $C_{2\ell}^{+2}$, or even beachballs, for all $\ell\geq2$
\cite[Theorem~4.3]{GLLV}. The one-apex even-cycle case is the
commonness theorem for even wheels. More recently, Kr\'a\v l, Krnc,
and Lamaison proved that $T^{+k}$ is common for every tree $T$ when
$1\leq k\leq5$ \cite[Theorem~1]{KKL}.
The results below remove the restriction on the apex number for
trees and establish the entire even-cycle family. Thus the tree
result extends the known range beyond five apices, while the new
range for even cycles is three or more apices.

\subsection{Main results}

Throughout, $P_j$ is the path with $j$ edges, and a tree is finite,
connected, and nonempty. In particular, $P_2$ has three vertices.

\begin{theorem}[Tree apices]\label{thm:trees}
For every tree $T$ and every integer $k\geq1$, the graph $T^{+k}$
is common. More precisely, if $T$ has $n\geq2$ vertices, then
\begin{equation}\label{eq:tree-two-color}
 M(T^{+k},W)\geq
 \frac{M(K_3,W)^{k(n-1)}}{M(P_2,W)^{(k-1)(n-2)}}
\end{equation}
for every graphon $W$. The denominator is positive, and equality
holds at $W\equiv1/2$.
\end{theorem}

The entropy argument actually proves a one-color inequality.

\begin{theorem}\label{thm:tree-one-color}
For a tree $T$ on $n\geq2$ vertices, an integer $k\geq1$, and a
graphon $W$,
\begin{equation}\label{eq:tree-one-color}
 t(T^{+k},W)t(K_2,W)^{n+k-3}
 t(P_2,W)^{(k-1)(n-2)}\geq t(K_3,W)^{k(n-1)}.
\end{equation}
Factors with exponent zero are omitted. Every constant graphon
attains equality.
\end{theorem}

To see why Theorem~\ref{thm:tree-one-color} implies commonness,
apply it to the graphon consisting of copies of $W$ and $1-W$ on
two blocks of measure $1/2$, with zero between the blocks. For a
connected graph $F$, its density becomes $2^{-|V(F)|}M(F,W)$.
The vertex counts make these powers of two cancel, and
$M(K_2,W)=1$. Writing $\sigma=M(P_2,W)$ and $\tau=M(K_3,W)$,
we therefore obtain
\[
 M(T^{+k},W)\geq
 \tau^k\bigl[\tau(\tau/\sigma)^{k-1}\bigr]^{n-2}
 \geq 2^{2-(k+1)n}=2^{1-e(T^{+k})}.
\]
The second inequality follows from Goodman's identity
$\tau=\tfrac32\sigma-\tfrac12$ and $\sigma\geq\tfrac12$, which
give $\tau\geq\tfrac14$ and $\tau/\sigma\geq\tfrac12$
\cite{Goodman}. This proves commonness for $n\geq2$; the
one-vertex tree gives a star, whose commonness follows from
convexity. The details appear in Section~\ref{sec:trees}.

\begin{theorem}[Even-cycle apices]\label{thm:cycles}
For every even $n\geq4$, every integer $k\geq1$, and every graphon $W$,
\begin{equation}\label{eq:cycle-common}
 M(C_n^{+k},W)\geq2^{1-n(k+1)}.
\end{equation}
Equality holds if and only if $W=1/2$ almost everywhere.
If $k\geq3$, then
\begin{equation}\label{eq:cycle-relative}
 M(C_n^{+k},W)\geq2^{-nk}M(C_n,W)\geq2^{1-n(k+1)}.
\end{equation}
\end{theorem}

The edge counts are $e(T^{+k})=(k+1)n-1$ and
$e(C_n^{+k})=n(k+1)$, so the displayed commonness bounds have the
required random-coloring normalization.

\subsection{Two proof mechanisms}

The tree argument uses relative entropy to turn a branching law
into a lower bound on a homomorphism density. This approach to
Sidorenko-type inequalities originates in Li and Szegedy's
logarithmic calculus \cite{LiSzegedy} and Szegedy's systematic
information-theoretic framework \cite{SzegedyEntropy}. The
tree-arrangeability work of Kim, Choongbum Lee, and Joonkyung Lee
\cite{KLL}, and the branching random walks of Conlon, Kim,
Choongbum Lee, and Joonkyung Lee \cite{CKLL}, develop the relevant
tree-based constructions. Higher tree decompositions were studied
by the same four authors \cite{CKLLHigher}.
Joonkyung Lee further applied entropy to graphs assembled along
tree decompositions in the setting of locally dense graphs
\cite{LeeLocal}.
Our local distribution is the normalized triangle density on a
finite weighted host. Conditionally independent extensions of its
common edge produce the apex set. Gibbs' inequality controls the
entropy cost of these extensions, and a branching construction
spreads the resulting edge law over the entire tree.

For cycles, we use the positivity principle of Razborov's flag
algebra method \cite{Razborov}, which also underlies the
computer-assisted part of \cite{GLLV}. We give it in a signed
homomorphism-density basis: a nonnegative rooted quadratic form
expands into a finite linear combination of signed graph densities.
Four exact identities on six vertices provide the needed comparison
at the fourth moment. The main analytic step converts this
comparison into a statement that retains the larger of the two
color moments. Convexity then raises the cycle length without
losing the distinction between the colors.

Section~\ref{sec:prelim} gives the approximation tools.
Section~\ref{sec:trees} develops the entropy argument, and
Section~\ref{sec:cycles} proves the spectral and moment comparisons.
The precise certificate data and verification are described in
Appendix~\ref{app:certificates}. The results apply to two classes
of connected Sidorenko graphs; they do not establish
Conjecture~\ref{conj:apices} for arbitrary Sidorenko bases.

\section{Density and approximation tools}\label{sec:prelim}

We use $t(F,L)$ for any bounded symmetric real kernel $L$ on a
probability space, with the same integral definition as above.
Vertex relabeling and the addition of isolated vertices leave the
density unchanged. Fubini's theorem also gives
\[
 t(F_1\sqcup F_2,L)=t(F_1,L)t(F_2,L).
\]
All graphs whose densities occur are simple. Thus changing $L$ on
a null set leaves their densities unchanged: the exceptional set
for any edge has zero measure under its two distinct vertex
coordinates.

\begin{lemma}[Continuity]\label{lem:continuity}
If $F$ has $e\geq1$ edges and $|L|,|L'|\leq B$ for $B\geq1$, then
\[
 |t(F,L)-t(F,L')|\leq eB^{e-1}\|L-L'\|_1.
\]
Every graphon is the $L^1$ limit of symmetric step graphons on
finite equal-size dyadic partitions. Cell averaging preserves its
range and its integral.
\end{lemma}
\begin{proof}
Telescope the difference of the two edge products. Each summand
has one factor $L-L'$ and $e-1$ factors bounded by $B$. Its absolute
integral is at most $B^{e-1}\|L-L'\|_1$, proving the estimate.
For the approximation, let $Q_jL$ be the average of $L$ on each
product of dyadic intervals of length $2^{-j}$. This operation
contracts $L^1$ and fixes every coarser dyadic step function.
Such step functions are dense in $L^1$: approximate by simple
functions and approximate their measurable sets in measure by
finite unions of dyadic rectangles. If $G$ is such an approximation,
then for all sufficiently large $j$,
\[
 \|Q_jL-L\|_1\leq\|Q_j(L-G)\|_1+\|G-L\|_1\leq2\|G-L\|_1.
\]
This proves convergence. Symmetry, bounds, and the integral are
preserved on each cell.
\end{proof}

On a finite graph $G$, we count all labeled homomorphisms and write
$t(F,G)=\hom(F,G)/|V(G)|^{|V(F)|}$. Hosts may have loops when
explicitly stated; source graphs remain simple.

\begin{lemma}[Approximation by finite graphs]\label{lem:finite-approx}
Given a graphon $W$ and finitely many simple graphs $F_1,\ldots,F_s$,
there are finite simple graphs $G_j$ such that
$t(F_i,G_j)\to t(F_i,W)$ for every $i$.
\end{lemma}
\begin{proof}
Sample independent uniform points $X_1,\ldots,X_N$ and, conditional
on these points, insert each edge $ab$ independently with probability
$W(X_a,X_b)$. For a graph $F$ with $v$ vertices, let $Y_F$ be the
average homomorphism indicator over injective maps into this host.
Then $\E Y_F=t(F,W)$. Indicators for maps with disjoint images are
independent. For any fixed image, the proportion of injective maps
meeting it is at most $v^2/N$, so $\operatorname{Var}(Y_F)\leq v^2/N$.
The difference between $Y_F$ and the ordinary homomorphism density
is at most $\binom v2/N$, by the collision bound. Thus each density
converges in probability to its graphon value. A union bound gives
simultaneous convergence for the finite list, from which a
deterministic sequence of realizations can be chosen.
\end{proof}

We shall repeatedly use the probability-space moment inequality
\begin{equation}\label{eq:moments}
 \E X^q\geq(\E X^p)^{q/p}
 \qquad(X\geq0,\ 1\leq p\leq q),
\end{equation}
which follows from Jensen. All real powers below have nonnegative
bases and positive exponents, unless an exponent-zero factor is
explicitly omitted.

\section{Tree apices through relative entropy}\label{sec:trees}

\subsection{Relative entropy and a branching law}

For a probability law $p$ on a finite set and a nonnegative
reference weight $\rho$ that is positive on the support of $p$,
define
\begin{equation}\label{eq:relative-entropy}
 \mathcal H_\rho(p)=\sum_{a:p(a)>0}
                    p(a)\log\frac{\rho(a)}{p(a)}.
\end{equation}
All logarithms are natural. If $Z_\rho=\sum_a\rho(a)$, then
\begin{equation}\label{eq:gibbs}
 \mathcal H_\rho(p)
 =\log Z_\rho-D_{\mathrm{KL}}(p\|\rho/Z_\rho)
 \leq\log Z_\rho.
\end{equation}
Here the Kullback--Leibler divergence is
\[
 D_{\mathrm{KL}}(p\|q)=\sum_{a:p(a)>0}p(a)\log\frac{p(a)}{q(a)}.
\]
In particular, $\mathcal H_\rho(p)=-D_{\mathrm{KL}}(p\|\rho)$
when $\rho$ is a probability law.
For any probability law $q$ positive on the support of $p$,
the inequality $\log u\leq u-1$ gives
\[
 \sum_{a:p(a)>0}p(a)\log\frac{q(a)}{p(a)}
 \leq\sum_{a:p(a)>0}q(a)-1\leq0.
\]
This proves Gibbs' inequality and \eqref{eq:gibbs}.
Taking $\rho\equiv1$ recovers Shannon entropy and its support
bound. More generally, if $\rho$ is the weight of a vertex map
in a homomorphism density, then $Z_\rho$ is precisely that
density. This choice absorbs the normalization into the entropy
calculation.

The branching construction below is the conditional version of
the random walk used in \cite[Section~2]{CKLL}, with reference
weights attached to the same tree.

\begin{lemma}[Tree extension]\label{lem:tree-extension}
Let $S,V$ be finite sets, and let $Q(s,x,y)$ be a probability
law on $S\times V\times V$ that is symmetric in $x,y$. Put
$\mu(s,x)=\sum_yQ(s,x,y)$ and
$\kappa_s(x,y)=Q(s,x,y)/\mu(s,x)$ whenever $\mu(s,x)>0$.
Let $\rho_S(s)$, $\nu_s(x)$, and a symmetric $\eta(x,y)$ be
nonnegative reference factors such that
\[
 b(s,x,y)=\rho_S(s)\nu_s(x)\nu_s(y)\eta(x,y)>0
 \qquad\text{whenever }Q(s,x,y)>0.
\]
For every tree $T$ on $n\geq2$ vertices, there is an extension
$P_T$ of $Q$ with the original law on every edge and with
\begin{equation}\label{eq:extension-entropy}
 \mathcal H_{\rho_T}(P_T)=\mathcal H_b(Q)+(n-2)g,
\end{equation}
where
\begin{align*}
 \rho_T(s,\mathbf x)
   &=\rho_S(s)\prod_{v\in V(T)}\nu_s(x_v)
                    \prod_{uv\in E(T)}\eta(x_u,x_v),\\
 g&=\sum_{Q(s,x,y)>0}Q(s,x,y)
                 \log\frac{\nu_s(y)\eta(x,y)}{\kappa_s(x,y)}.
\end{align*}
Moreover, $\rho_T$ is positive on the support of $P_T$.
\end{lemma}
\begin{proof}
Start from an edge $uv$ and sample $(s,x_u,x_v)$ with law $Q$.
Orient the remaining edges away from this edge and sample each
new vertex from $\kappa_s$ at its parent. Rows with
$\mu(s,x)=0$ may be assigned arbitrary probability laws: such
states are never visited. Symmetry gives
\[
 \sum_x\mu(s,x)\kappa_s(x,y)
   =\sum_xQ(s,x,y)=\mu(s,y).
\]
Induction therefore gives marginal $\mu$ at each vertex and
law $Q$ on each edge. The reference factors are positive on
every sampled edge, proving the support assertion.

On this support, $\rho_T/P_T$ is the product of $b/Q$ on the
initial edge and $\nu_s(y)\eta(x,y)/\kappa_s(x,y)$ for each of
the $n-2$ subsequently added vertices. Taking logarithms and
expectations proves \eqref{eq:extension-entropy}.
\end{proof}

\subsection{The weighted triangle law}

A finite weighted host consists of a finite set $V$, weights
$w_x\geq0$ with $\sum_xw_x=1$, and a symmetric kernel
$M:V^2\to[0,1]$. Diagonal entries are allowed. Its density is
\[
 t(F,M)=\sum_{\varphi:V(F)\to V}
         \prod_{v\in V(F)}w_{\varphi(v)}
         \prod_{uv\in E(F)}M(\varphi(u),\varphi(v)).
\]
Define
\begin{align*}
 d(x)&=\sum_yw_yM(x,y),&
 r(x,y)&=\sum_zw_zM(x,z)M(y,z),\\
 a(x)&=\sum_yw_yM(x,y)r(x,y),&
 E&=t(K_2,M),\quad D=t(P_2,M),\quad R=t(K_3,M).
\end{align*}
Thus $E=\sum_xw_xd(x)$, $D=\sum_xw_xd(x)^2$, and
$R=\sum_xw_xa(x)$. Assume $R>0$. Sample the triangle law
\[
 p_\triangle(x,y,z)
   =\frac{w_xw_yw_zM(x,y)M(x,z)M(y,z)}{R}.
\]
It is invariant under coordinate permutations, and its pair and
vertex marginals are
\[
 p_2(x,y)=\frac{w_xw_yM(x,y)r(x,y)}{R},\qquad
 p_1(x)=\frac{w_xa(x)}{R}.
\]
Put
\[
 h=\sum_{x,y}p_2(x,y)\log r(x,y),\qquad
 \alpha=\sum_xp_1(x)\log a(x).
\]
Terms with zero probability are omitted here and below.
Since $0\leq a(x)\leq d(x)^2\leq d(x)$, the assumption
$R>0$ implies $D,E>0$.

\begin{lemma}\label{lem:triangle}
The weighted triangle law satisfies
\begin{equation}\label{eq:triangle-estimates}
 h\geq\log(R/E),\qquad \alpha-2h\leq\log(D/R).
\end{equation}
\end{lemma}
\begin{proof}
Apply \eqref{eq:gibbs} to $p_2$ with reference
$\rho(x,y)=w_xw_yM(x,y)$, whose total mass is $E$.
On the support of $p_2$, $\rho/p_2=R/r(x,y)$, so
$\log R-h\leq\log E$.

For the second estimate, write
$\ell=\sum_xp_1(x)\log d(x)$. The comparison law
\[
 q(x,y)=p_1(x)\frac{w_yM(x,y)}{d(x)}
\]
is set to zero when $d(x)=0$. It has total mass one and is
positive wherever $p_2$ is positive. Therefore
\[
 \alpha-\ell-h=\mathcal H_q(p_2)\leq0.
\]
Applying \eqref{eq:gibbs} to $p_1$ with reference
$\widetilde\rho(x)=w_xd(x)^2/R$, of total mass $D/R$, gives
\[
 2\ell-\alpha=\mathcal H_{\widetilde\rho}(p_1)
                    \leq\log(D/R).
\]
Adding twice the first inequality to the second proves
$\alpha-2h\leq\log(D/R)$.
\end{proof}

\subsection{The book law and the page estimate}

Retain the pair $(X,Y)$ with law $p_2$. Conditional on it,
sample $Z_1,\ldots,Z_k$ independently with probabilities
$w_zM(X,z)M(Y,z)/r(X,Y)$. Write $s=(z_1,\ldots,z_k)$ and set
\[
 \rho_S(s)=\prod_{i=1}^kw_{z_i},\qquad
 \nu_s(x)=w_x\prod_{i=1}^kM(x,z_i),\qquad \eta(x,y)=M(x,y).
\]
The resulting book law is symmetric in $x,y$ and has the form
\begin{equation}\label{eq:book-law}
 Q(s,x,y)=\frac{b(s,x,y)}{R\,r(x,y)^{k-1}},\qquad
 b(s,x,y)=\rho_S(s)\nu_s(x)\nu_s(y)M(x,y),
\end{equation}
with $Q=0$ when $r(x,y)=0$. Every triple $(X,Y,Z_i)$ has
law $p_\triangle$. Let $\mu(s,x)=\sum_yQ(s,x,y)$ and let $g$
be the conditional relative entropy in
Lemma~\ref{lem:tree-extension}, for these reference factors.

\begin{lemma}\label{lem:book}
For the book law,
\begin{align}
 \mathcal H_b(Q)&\geq k\log R-(k-1)\log E,
                                      \label{eq:book-total}\\
 g&\geq k\log R-\log E-(k-1)\log D.
                                      \label{eq:book-step}
\end{align}
\end{lemma}
\begin{proof}
Since $b/Q=Rr(x,y)^{k-1}$ on the support,
\[
 \mathcal H_b(Q)=\log R+(k-1)h.
\]
The first estimate follows from \eqref{eq:triangle-estimates}.

For the second, compare $\mu$ with the probability law
\[
 \lambda(s,x)=p_1(x)\prod_{i=1}^k\frac{p_2(x,z_i)}{p_1(x)},
\]
set to zero when $p_1(x)=0$. This law has independently sampled
pages conditional on $X=x$. It is positive on the support of
$\mu$, because each $(X,Z_i)$ marginal of the book law is $p_2$.
Using the explicit marginals, on the support of $Q$ we have
\[
 \frac{\nu_s(y)M(x,y)}{\kappa_s(x,y)}
 =\frac{\mu(s,x)}{\lambda(s,x)}
   \frac{r(x,y)^{k-1}\prod_{i=1}^kr(x,z_i)}{a(x)^{k-1}}.
\]
Taking logarithms and expectations gives the identity
\begin{equation}\label{eq:pages-kl}
 g=h-(k-1)(\alpha-2h)+D_{\mathrm{KL}}(\mu\|\lambda).
\end{equation}
Indeed, each $\log r(X,Z_i)$ has expectation $h$, as does
$\log r(X,Y)$, while $\log a(X)$ has expectation $\alpha$.
Nonnegativity of the divergence and
\eqref{eq:triangle-estimates} now yield
\[
 g\geq\log(R/E)-(k-1)\log(D/R),
\]
which is \eqref{eq:book-step}.
\end{proof}

\subsection{The density bound and transfer to graphons}

\begin{proposition}\label{prop:tree-finite}
For every finite weighted host, every tree $T$ on $n\geq2$
vertices, and every $k\geq1$,
\begin{equation}\label{eq:tree-finite}
 t(T^{+k},M)E^{n+k-3}D^{(k-1)(n-2)}\geq R^{k(n-1)}.
\end{equation}
\end{proposition}
\begin{proof}
If $R=0$, the right side vanishes. Otherwise extend the book law
along $T$ by Lemma~\ref{lem:tree-extension}. Its reference weight
is
\[
 \rho_T(s,\mathbf x)
 =\left(\prod_{i=1}^kw_{z_i}\right)
  \left(\prod_{v\in V(T)}w_{x_v}\prod_{i=1}^kM(x_v,z_i)\right)
  \prod_{uv\in E(T)}M(x_u,x_v).
\]
Its total mass is exactly $t(T^{+k},M)$, and it is positive on
the support of the extension. Gibbs' inequality, followed by
\eqref{eq:extension-entropy} and Lemma~\ref{lem:book}, gives
\begin{align*}
 \log t(T^{+k},M)
 &\geq\mathcal H_{\rho_T}(P_T)
       =\mathcal H_b(Q)+(n-2)g\\
 &\geq k(n-1)\log R-(n+k-3)\log E
                       -(k-1)(n-2)\log D.
\end{align*}
Exponentiating proves \eqref{eq:tree-finite}.
\end{proof}

\begin{proof}[Proof of Theorem~\ref{thm:tree-one-color}]
Every step graphon is a finite weighted host: its vertex weights
are the measures of its cells, and its kernel entries are the
cell values. Proposition~\ref{prop:tree-finite} therefore gives
\eqref{eq:tree-one-color} for step graphons.
Lemma~\ref{lem:continuity} extends it to every graphon, since
both sides are continuous polynomials in four fixed densities.
For $W\equiv p$, put $a=n+k-3$, $b=(k-1)(n-2)$, and
$c=k(n-1)$. The total edge exponent on the left is
$((k+1)n-1)+a+2b=3c$, equal to that on the right.
This proves the equality assertion, including $p=0$.
\end{proof}

\subsection{Combining the colors}

Put $a=n+k-3$, $b=(k-1)(n-2)$, and $c=k(n-1)$.
The vertex counts satisfy
\begin{equation}\label{eq:vertex-balance}
 (n+k)+2a+3b=3c.
\end{equation}
Let $W\sqcup(1-W)$ be the graphon with two blocks of measure $1/2$,
carrying rescaled copies of $W$ and $1-W$, and zero between the
blocks. For every connected graph $F$,
\begin{equation}\label{eq:disjoint-colors}
 t(F,W\sqcup(1-W))=2^{-|V(F)|}M(F,W).
\end{equation}
Indeed, a nonzero contribution from a connected graph with an edge
must lie in one block. Each block contributes its density times
$2^{-|V(F)|}$; the single-vertex case is immediate.
Apply Theorem~\ref{thm:tree-one-color} to this graphon and use
\eqref{eq:vertex-balance}. Since $M(K_2,W)=1$, the powers of two
cancel to give
\begin{equation}\label{eq:tree-polynomial-colors}
 M(T^{+k},W)M(P_2,W)^{(k-1)(n-2)}
 \geq M(K_3,W)^{k(n-1)}.
\end{equation}

For completeness, put $\sigma=M(P_2,W)$ and $\tau=M(K_3,W)$.
Goodman's identity \cite{Goodman} takes the form
\begin{equation}\label{eq:goodman}
 \sigma\geq\tfrac12,\qquad
 \tau=\tfrac32\sigma-\tfrac12\geq\tfrac14,
 \qquad \tau/\sigma\geq\tfrac12.
\end{equation}
To see this directly, if $d_W(x)=\int W(x,y)\dd y$, then
\[
 \sigma=\int\bigl(d_W^2+(1-d_W)^2\bigr)
         =\tfrac12+2\int(d_W-\tfrac12)^2.
\]
Integrating
$abc+(1-a)(1-b)(1-c)=1-a-b-c+ab+bc+ca$
around a triangle yields the identity relating $\tau$ and $\sigma$.
The other bounds follow immediately.

\begin{proof}[Proof of Theorem~\ref{thm:trees}]
Division in \eqref{eq:tree-polynomial-colors} is valid by
\eqref{eq:goodman}, proving \eqref{eq:tree-two-color}. Moreover,
\begin{align*}
 M(T^{+k},W)
 &\geq\tau^k\left[\tau(\tau/\sigma)^{k-1}\right]^{n-2}\\
 &\geq 4^{-k}\left[\tfrac14\,2^{-(k-1)}\right]^{n-2}
   =2^{2-(k+1)n}.
\end{align*}
This is the commonness bound when $n\geq2$. If $T$ has one vertex,
then $T^{+k}=K_{1,k}$, and pointwise convexity gives
\[
 M(K_{1,k},W)=\int\bigl(d_W(x)^k+(1-d_W(x))^k\bigr)\dd x
             \geq2^{1-k}.
\]
Finally, $W\equiv1/2$ has $\sigma=1/2$ and $\tau=1/4$, and
attains the quantitative bound.
\end{proof}

\section{Even-cycle apices through fourth moments}\label{sec:cycles}

The proof starts on a finite uniform probability space with $d$
points. Lemma~\ref{lem:continuity} will transfer the final bounds
to graphons. Given $W$, use its signed kernel $U=2W-1$ and write
\[
 S_\varepsilon=1+\varepsilon U,\qquad \varepsilon\in\{-1,1\},
\]
where the sign is sampled uniformly. Averages over host coordinates
use the uniform probability measure. Thus $S_+=2W$ and
$S_-=2(1-W)$, and
\begin{equation}\label{eq:signed-average}
 \E_\varepsilon t(F,S_\varepsilon)=2^{e(F)-1}M(F,W).
\end{equation}
Expanding its left side keeps the even edge subsets of $F$;
multiplication by $\varepsilon$ instead keeps the odd subsets.

\subsection{Base-cycle traces}

Use the inner product $\ip xy_d=d^{-1}\sum_i x_i y_i$, so the
constant vector $\one$ has norm one. The operator associated with
a kernel $L$ is $T_L=L/d$. Expanding the trace gives
\begin{equation}\label{eq:cycle-trace}
 t(C_n,L)=\Tr T_L^n\qquad(n\geq3).
\end{equation}
If $B$ is real symmetric, $g$ is
a unit vector, and $n\geq4$ is even, diagonalization gives
\begin{equation}\label{eq:spectral}
 \Tr B^n\geq\|Bg\|^n\geq|\ip g{Bg}|^n,
 \qquad \Tr B^n\leq(\Tr B^4)^{n/4}.
\end{equation}
For the last inequality, apply
$\sum_i x_i^s\leq(\sum_i x_i)^s$ to the nonnegative fourth
powers of the eigenvalues, with $s=n/4$.

The signed densities needed below satisfy
\begin{align}
 0&\leq t(K_2,U)^2\leq t(P_2,U)\leq1,
 \qquad 0\leq t(C_4,U)\leq1,\notag\\
 |t(P_3,U)|&\leq t(P_2,U)t(C_4,U)^{1/4}.
 \label{eq:scalar-bounds}
\end{align}
Indeed, $t(P_2,U)=\|T_U\one\|^2$, whereas
$t(K_2,U)=\ip\one{T_U\one}_d$.
Also $\Tr T_U^2=\E U^2\leq1$ implies
$\Tr T_U^4\leq1$ and
$\|T_U\|_\op\leq t(C_4,U)^{1/4}$.
Apply this operator bound to
$t(P_3,U)=\ip{T_U\one}{T_U(T_U\one)}_d$.

Expanding around a four-cycle gives
\begin{align}
 8M(C_4,W)&=1+2t(K_2,U)^2+4t(P_2,U)+t(C_4,U),\notag\\
 t(C_4,S_\varepsilon)
 &=8M(C_4,W)+4\varepsilon\bigl(t(K_2,U)+t(P_3,U)\bigr).
 \label{eq:fourth-colors}
\end{align}
The six two-edge subsets consist of two matchings and four
adjacent pairs, while the four three-edge subsets are paths.
This explains all coefficients. The Rayleigh value of
$T_{S_\varepsilon}$ at $\one$ is $1+\varepsilon t(K_2,U)$.
Thus, for even $n\geq4$,
\begin{equation}\label{eq:even-base}
 2^{n-1}M(C_n,W)
 \geq\frac{(1+t(K_2,U))^n+(1-t(K_2,U))^n}{2}
 \geq1+\binom n2t(K_2,U)^2\geq1.
\end{equation}

\subsection{Conditioning on the apices}

For $s\geq1$, sample the sign and independent apex coordinates
$z_1,\ldots,z_s$. Define
\[
 h_s(x)=\prod_{i=1}^s S_\varepsilon(z_i,x),\qquad
 D_s=\E_xh_s(x),\qquad
 \Pi_s=\E_xh_s(x)\bigl(\E_yS_\varepsilon(x,y)h_s(y)\bigr)^2.
\]
Let $V_s=(\Pi_s/D_s)^{1/2}$ when $D_s>0$, and let $V_s=0$
otherwise. If $D_s=0$, nonnegativity forces $h_s=0$, and hence
$\Pi_s=0$. In particular, $D_sV_s^2=\Pi_s$ in all cases.
Since $0\leq S_\varepsilon\leq2$, we have $\Pi_s\leq4D_s^3$ and
$0\leq V_s\leq2D_s\leq2^{s+1}$. Thus all moments below are finite.
Unless subscripts specify otherwise, expectations in this part of
the proof include the sign and all apex coordinates.

\begin{lemma}\label{lem:conditional}
For every even $n\geq4$,
\begin{equation}\label{eq:conditional-cycle}
 2^{n(s+1)-1}M(C_n^{+s},W)\geq\E V_s^n.
\end{equation}
If $\E D_s^2>0$, then
\begin{equation}\label{eq:fourth-support}
 V_s^4\geq2\Pi_s-D_s^2,
 \qquad \E V_s^4\geq\frac{(\E\Pi_s)^2}{\E D_s^2}.
\end{equation}
\end{lemma}
\begin{proof}
Put $B_s=\diag(\sqrt{h_s})T_{S_\varepsilon}\diag(\sqrt{h_s})$.
Expanding its trace puts a factor $h_s$ at every cycle vertex,
so $\E\Tr B_s^n=2^{n(s+1)-1}M(C_n^{+s},W)$.
When $D_s>0$, the vector
$g_s=\sqrt{h_s/D_s}$ has norm one and
$\|B_sg_s\|^2=\Pi_s/D_s=V_s^2$. Equation~\eqref{eq:spectral}
therefore proves \eqref{eq:conditional-cycle}; the zero case has
$B_s=0$.
For positive $D_s$, the first inequality in
\eqref{eq:fourth-support} is the identity
\[
 V_s^4-2\Pi_s+D_s^2=(\Pi_s/D_s-D_s)^2.
\]
It remains valid in the zero case. Finally, Cauchy--Schwarz applied
to $\E\Pi_s=\E D_sV_s^2$ proves the second assertion.
\end{proof}

These conditional quantities have simple density interpretations:
\begin{equation}\label{eq:conditional-densities}
 \E D_s^2=2^{2s-1}M(K_{2,s},W),\qquad
 \E\Pi_s=2^{3s+1}M(K_{1,2,s},W).
\end{equation}
Here $K_{1,2,s}$ is the complete tripartite graph: the $s$ apices
meet three further vertices, and one of those vertices meets the
other two. In particular, integrating the apices individually and
applying \eqref{eq:moments} gives
\begin{align}
 \E D_2^2&=8M(C_4,W),\notag\\
 \E D_3^2
 &=\E_{\varepsilon,x,y}
     \bigl(\E_zS_\varepsilon(x,z)S_\varepsilon(y,z)\bigr)^3
 \geq\bigl(8M(C_4,W)\bigr)^{3/2}.
 \label{eq:codegree}
\end{align}

For three apices, introduce
\begin{equation}\label{eq:majority}
 \maj=U(z_1,z_2)+U(z_1,z_3)+U(z_2,z_3)
                  -U(z_1,z_2)U(z_1,z_3)U(z_2,z_3)
\end{equation}
Its mean and its role in the conditional fourth-moment bound are
recorded as
\[
 \mu=\E\maj=3t(K_2,U)-t(K_3,U),\qquad F=2\Pi_3-D_3^2.
\]
The polynomial $x+y+z-xyz$ is multiaffine on $[-1,1]^3$ and has
values $\pm2$ at its corners. Therefore
\begin{equation}\label{eq:majority-bounds}
 |\maj|\leq2.
\end{equation}

The following is the entire certificate input to the proof.

\begin{lemma}[Six-vertex inequalities]\label{lem:cert-input}
For every $W$ and its signed kernel $U=2W-1$,
\begin{equation}\label{eq:mean-input}
 M(K_{1,2,2},W)\geq\frac{M(C_4,W)}{16},\qquad
 M(K_{1,2,3},W)\geq\frac{M(K_{2,3},W)}{32}.
\end{equation}
Moreover, the conditional quantity $F$ satisfies
\begin{align}
 \E(3+\varepsilon-\maj)F
 &\geq24M(C_4,W)+13t(K_2,U)+t(K_3,U)\notag\\
 &\qquad+12t(K_2,U)^2-4t(K_2,U)t(K_3,U)+12t(P_3,U),
 \label{eq:G}\\
 \E(6+\varepsilon+\maj)F
 &\geq48M(C_4,W)+19t(K_2,U)+23t(K_3,U)\notag\\
 &\qquad-12t(K_2,U)^2+4t(K_2,U)t(K_3,U)+24t(P_3,U).
 \label{eq:H}
\end{align}
\end{lemma}
\begin{proof}
Appendix~\ref{app:certificates} specifies four rational sums of
nonnegative rooted quadratic forms. Their evaluated targets are
respectively $\E\Pi_2-\E D_2^2$, $\E\Pi_3-\E D_3^2$, and
the left sides minus the right sides of \eqref{eq:G} and
\eqref{eq:H}. The interpretation
\eqref{eq:conditional-densities} gives \eqref{eq:mean-input}.
The coefficient identities and positive factorizations are checked
exactly there. In particular, the inequalities impose no sign
condition on $t(K_2,U)$.
\end{proof}

Combining \eqref{eq:fourth-support}, \eqref{eq:codegree}, and
\eqref{eq:mean-input} gives
\begin{equation}\label{eq:unweighted}
 \E V_2^4\geq8M(C_4,W),\qquad
 \E V_3^4\geq\E F\geq\bigl(8M(C_4,W)\bigr)^{3/2}
                          \geq8M(C_4,W)\geq1.
\end{equation}
The variable $F$ itself need not be nonnegative. Its usefulness
comes from lying below $V_3^4$ pointwise.

\subsection{Retaining the larger color moment}

The average bound in \eqref{eq:unweighted} does not by itself
compare higher moments separately for the two colors. We need a
probability weight that recovers the larger fourth color trace.
The following scalar inequality will be applied to the densities
in \eqref{eq:fourth-colors}.

\begin{lemma}[Scalar estimate]\label{lem:scalar-comparison}
If $b\geq0$, $0\leq r\leq1$, and $a\geq1+4b+r^4$, then
\begin{equation}\label{eq:scalar-comparison}
 a^{3/2}\geq a+4br.
\end{equation}
\end{lemma}
\begin{proof}
If $r\leq1/2$, then $a\geq1$ gives
\[
 a^{3/2}-a=\frac{a(a-1)}{\sqrt a+1}
             \geq\frac{a-1}{2}\geq2b\geq4br.
\]
If $r\geq1/2$, the function $y^{3/2}-y$ is increasing on
$[1,\infty)$, so it suffices to take $a=1+4b+r^4$.
For $x=\sqrt{1+4b+r^4}$ the required difference becomes
$x^3-(1+r)x^2+r+r^5$. The identity
\begin{align*}
 x^3-(1+r)x^2+r+r^5
 &=(x-\tfrac23(1+r))^2(x+\tfrac13(1+r))\\
 &\qquad+\tfrac1{27}\bigl(27r^5-4(1+r)^3+27r\bigr)
\end{align*}
reduces its nonnegativity to that of the last polynomial.
With $t=r-1/2\geq0$, this polynomial equals
\[
 27t^5+\frac{135}{2}t^4+\frac{127}{2}t^3
       +\frac{63}{4}t^2+\frac{135}{16}t+\frac{27}{32},
\]
whose coefficients are positive.
\end{proof}

\begin{proposition}[Weighted comparison]\label{prop:weighted}
There is a weight $\omega$, depending on the sampled sign and
three apices, such that
\begin{equation}\label{eq:weighted}
 0\leq\omega\leq2,\qquad \E\omega=1,
 \qquad \E\omega V_3^4
       \geq16\max\{t(C_4,W),t(C_4,1-W)\}.
\end{equation}
\end{proposition}
\begin{proof}
Complement first so that $t(K_2,U)\geq0$. By
\eqref{eq:fourth-colors}, Lemma~\ref{lem:scalar-comparison}
applies with $a=8M(C_4,W)$, $b=t(P_2,U)$, and
$r=t(C_4,U)^{1/4}$.

If $t(K_2,U)+t(P_3,U)\leq0$, then the larger color density is
$t(C_4,1-W)$. The bound \eqref{eq:scalar-bounds} gives
\begin{align*}
 \E V_3^4
 &\geq\bigl(8M(C_4,W)\bigr)^{3/2}\\
 &\geq8M(C_4,W)+4t(P_2,U)t(C_4,U)^{1/4}\\
 &\geq8M(C_4,W)-4\bigl(t(K_2,U)+t(P_3,U)\bigr)
   =16t(C_4,1-W).
\end{align*}
Thus $\omega=1$ suffices in this case.

Otherwise $t(C_4,W)$ is larger. We will also use
\begin{equation}\label{eq:scalar-remainder}
 2\E F-8M(C_4,W)-4t(P_3,U)-1\geq0.
\end{equation}
Indeed, the negative-color Rayleigh bound gives
\begin{align*}
 8M(C_4,W)-4t(P_3,U)
 &=16t(C_4,1-W)+4t(K_2,U)\\
 &\geq(1-t(K_2,U))^4+4t(K_2,U)\geq1,
\end{align*}
using $(1-x)^4+4x-1=x^2((x-2)^2+2)$.
Together with $\E F\geq8M(C_4,W)$ from
\eqref{eq:unweighted}, this proves \eqref{eq:scalar-remainder}.
Choose
\begin{equation}\label{eq:weight-formula}
 \omega=
 \begin{cases}
 (3+2\mu+\varepsilon-\maj)/(3+\mu),&\mu\geq0,\\[2pt]
 (6-2\mu+\varepsilon+\maj)/(6-\mu),&\mu<0.
 \end{cases}
\end{equation}

Suppose first that $\mu\geq0$. By \eqref{eq:majority-bounds},
the first numerator in \eqref{eq:weight-formula}
lies between $2\mu$ and $6+2\mu$, and has mean $3+\mu$.
Thus $0\leq\omega\leq2$ and $\E\omega=1$.
Multiply the pointwise inequality $V_3^4\geq F$ by this
nonnegative weight. Substitution of \eqref{eq:G} gives
\begin{align*}
 (3+\mu)\bigl(\E\omega V_3^4-16t(C_4,W)\bigr)
 &\geq(3+2\mu)\E F+\E(\varepsilon-\maj)F
                                  -16(3+\mu)t(C_4,W)\\
 &\geq4t(K_2,U)
     +\mu\bigl(2\E F-8M(C_4,W)-4t(P_3,U)-1\bigr)\geq0.
\end{align*}
The second step substitutes \eqref{eq:G},
$\mu=3t(K_2,U)-t(K_3,U)$, and \eqref{eq:fourth-colors};
the last uses \eqref{eq:scalar-remainder}.

Finally, if $\mu<0$, the second numerator lies
between $3-2\mu$ and $9-2\mu$, within $[0,2(6-\mu)]$,
and has mean $6-\mu$. Its weight is again admissible. Now
\begin{align*}
 (6-\mu)\bigl(\E\omega V_3^4-16t(C_4,W)\bigr)
 &\geq(6-2\mu)\E F+\E(\varepsilon+\maj)F
                                  -16(6-\mu)t(C_4,W)\\
 &\geq24t(K_3,U)-8t(K_2,U)\\
 &\qquad-\mu\bigl(2\E F-8M(C_4,W)-4t(P_3,U)-1\bigr)
 \geq0.
\end{align*}
Here the second step uses \eqref{eq:H}, and $\mu<0$ implies
$t(K_3,U)>3t(K_2,U)\geq0$. The last step again follows from
\eqref{eq:scalar-remainder}. This proves all cases.
If a complementation was made, interchange the sign labels to
transport the constructed weight back to the original kernel.
\end{proof}

\subsection{Lifting the cycle length}

\begin{lemma}[Two-moment lifting]\label{lem:length-lifting}
Let $V\geq0$ be bounded on a probability space. Suppose
$a,b\geq0$ and a weight $\omega$ satisfy
\[
 0\leq\omega\leq2,\quad\E\omega=1,\quad
 \E V^4\geq(a+b)/2,\quad
 \E\omega V^4\geq\max(a,b).
\]
Then for every real $q\geq4$,
\begin{equation}\label{eq:length-lifting}
 \E V^q\geq\tfrac12(a^{q/4}+b^{q/4}).
\end{equation}
\end{lemma}
\begin{proof}
Interchanging $a,b$ if necessary, assume $a\geq b$, and set
$\ell=\E\omega V^4$ and $u=\E(2-\omega)V^4$.
Then $\ell\geq a$ and $\ell+u\geq a+b$, which imply
$\ell^s+u^s\geq a^s+b^s$ for $s\geq1$.
Indeed, if $u\geq b$, compare termwise. Otherwise write
$u=b-\delta$ with $0<\delta\leq b$. The sum condition gives
$\ell\geq a+\delta$, and
$(a+\delta)^s+(b-\delta)^s\geq a^s+b^s$:
the derivative of $(a+t)^s+(b-t)^s$ is nonnegative for
$0\leq t\leq\delta$.
Both $\omega$ and $2-\omega$ have mean one, so Jensen with
$s=q/4$ under the two weighted probability measures yields
\[
 \E V^q\geq\tfrac12(\ell^s+u^s)
             \geq\tfrac12(a^s+b^s).
\]
\end{proof}

Apply this lemma with $V=V_3$, $a=16t(C_4,W)$, and
$b=16t(C_4,1-W)$. Its hypotheses follow from
\eqref{eq:unweighted} and Proposition~\ref{prop:weighted}.
For every even $n\geq4$, \eqref{eq:conditional-cycle},
\eqref{eq:length-lifting}, and \eqref{eq:spectral} give
\begin{align}
 2^{4n-1}M(C_n^{+3},W)
 &\geq\E V_3^n\notag\\
 &\geq2^{n-1}\bigl(t(C_4,W)^{n/4}+t(C_4,1-W)^{n/4}\bigr)
 \notag\\
 &\geq2^{n-1}M(C_n,W).
 \label{eq:three-apices}
\end{align}

\subsection{Lifting the apex number and handling small numbers}

The apex-number argument applies directly on any probability
space. Consider a measure $\nu$ on two copies of the space of
cycle coordinates $x_0,\ldots,x_{n-1}$. On the first copy its
density is $\prod_{i=0}^{n-1}W(x_i,x_{i+1})$, with $x_n=x_0$;
on the second copy replace $W$ by $1-W$.
The total mass is $M(C_n,W)$. On the first copy put
\[
 \xi=\int\prod_{i=0}^{n-1}W(z,x_i)\dd z,
\]
and make the same replacement on the second.
Integrating the apex coordinates separately gives
$M(C_n^{+j},W)=\int\xi^j\dd\nu$ for every $j\geq1$.
Since $M(C_n,W)>0$ by \eqref{eq:even-base}, normalizing
$\nu$ and using Jensen gives
\begin{equation}\label{eq:apex-lifting}
 M(C_n^{+k},W)\geq
 \frac{M(C_n^{+s},W)^{k/s}}{M(C_n,W)^{k/s-1}}
 \qquad(k\geq s\geq1).
\end{equation}
Together with \eqref{eq:three-apices}, this proves
$M(C_n^{+k},W)\geq2^{-nk}M(C_n,W)$ for all $k\geq3$.

For two apices, \eqref{eq:conditional-cycle},
\eqref{eq:moments}, and \eqref{eq:unweighted} give
\begin{equation}\label{eq:two-apices}
 M(C_n^{+2},W)
 \geq2^{1-3n}\bigl(8M(C_4,W)\bigr)^{n/4}
 \geq2^{1-3n}.
\end{equation}
For one apex, \eqref{eq:conditional-densities} gives
\[
 \E D_1^2=2M(P_2,W)=1+t(P_2,U),\qquad
 \E\Pi_1=16M(K_4-e,W),
\]
where $K_4-e=K_{1,2,1}$ is the diamond.
Its even edge-subsets are the empty graph, two two-edge
matchings, eight adjacent pairs, one four-cycle, and four
triangles with a pendant edge. The signed density of the last
graph has absolute value at most
$\sqrt{t(P_2,U)t(C_4,U)}$, since Cauchy--Schwarz and $|U|\leq1$
give
\[
 \left|\E_{x,y}U(x,y)
   \bigl(\E_zU(x,z)\bigr)
   \bigl(\E_wU(x,w)U(y,w)\bigr)\right|
 \leq\sqrt{t(P_2,U)t(C_4,U)}.
\]
The expansion therefore yields
\begin{align}
 16M(K_4-e,W)
 &\geq1+2t(K_2,U)^2+8t(P_2,U)+t(C_4,U)\notag\\
 &\qquad-4\sqrt{t(P_2,U)t(C_4,U)}\label{eq:diamond}\\
 &\geq1+2t(P_2,U)+\tfrac13t(C_4,U).
 \label{eq:diamond-bound}
\end{align}
The last difference is at least
$2t(K_2,U)^2+
 (\sqrt{6t(P_2,U)}-\sqrt{2t(C_4,U)/3})^2$.
In particular, $16M(K_4-e,W)\geq2M(P_2,W)\geq1$.
Now \eqref{eq:fourth-support} and \eqref{eq:moments} imply
\begin{equation}\label{eq:one-apex}
 M(C_n^{+1},W)
 \geq2\left(\frac{M(K_4-e,W)^2}{2M(P_2,W)}\right)^{n/4}
 \geq2^{1-2n}.
\end{equation}
Thus every positive apex number is covered on finite uniform
spaces.

\subsection{Graphons and rigidity}

\begin{proposition}\label{prop:graphon-bounds}
For every graphon $W$, every even $n\geq4$, and every $k\geq1$,
the bounds \eqref{eq:one-apex} and \eqref{eq:two-apices} hold,
and $M(C_n^{+k},W)\geq2^{-nk}M(C_n,W)\geq2^{1-n(k+1)}$
for $k\geq3$.
\end{proposition}
\begin{proof}
Every quantity in these final bounds is a density of a fixed
simple graph. Lemma~\ref{lem:continuity} therefore makes it
continuous under $L^1$ approximation. Each dyadic step graphon
has precisely the densities of its finite uniform matrix,
including diagonal cell entries. Apply the finite inequalities
and pass to the limit. The denominator $2M(P_2,W)$ is at least
one by \eqref{eq:goodman}, and every power involved is continuous
on its nonnegative domain. The same approximation transfers the
signed-density bounds \eqref{eq:scalar-bounds},
\eqref{eq:fourth-colors}, \eqref{eq:even-base}, and
\eqref{eq:diamond-bound}, which will be used for equality.
\end{proof}

We finish by identifying the equality case. We use here the
standard spectral description of symmetric bounded kernels;
see \cite{Lovasz} for the graphon and kernel-operator framework.
For such a kernel $L$, its integral operator $T_L$ is compact,
self-adjoint, and Hilbert--Schmidt on $L^2([0,1])$.
For even $n=2\ell\geq4$, Fubini and the spectral theorem give
\begin{equation}\label{eq:infinite-trace}
 t(C_n,L)=\|T_L^\ell\|_\HS^2
          =\sum_i\lambda_i(T_L)^n.
\end{equation}
The sums converge because $T_L$ is Hilbert--Schmidt.
In particular,
\begin{equation}\label{eq:c-zero}
 t(C_4,U)=0\quad\Longleftrightarrow\quad
 U=0\text{ almost everywhere}.
\end{equation}
Indeed, a vanishing sum of fourth powers makes $T_U=0$ by the
compact self-adjoint spectral theorem, and
$\|T_U\|_\HS^2=\iint U^2$ then vanishes.

\begin{lemma}\label{lem:cycle-equality}
For every graphon and every even $n\geq4$,
$M(C_n,W)=2^{1-n}$ holds if and only if
$W=1/2$ almost everywhere.
\end{lemma}
\begin{proof}
The forward implication starts with the graphon version of
\eqref{eq:even-base}, which forces $t(K_2,U)=0$. The constant unit
function consequently has Rayleigh value one for both
$T_{S_+}$ and $T_{S_-}$. Each operator has a largest eigenvalue
at least one, so each even cycle trace is at least one.
Their average is $2^{n-1}M(C_n,W)=1$, hence both traces equal one.
Equation~\eqref{eq:infinite-trace} then forces each spectrum
to consist of a single eigenvalue one and otherwise zeros.
Equality in the Rayleigh bound shows that the constant function
spans this eigenspace. Therefore each operator is the orthogonal
projection onto the constants, whose kernel is one almost
everywhere. It follows that $U=0$ almost everywhere.
Conversely, $W\equiv1/2$ gives $M(C_n,W)=2^{1-n}$ directly.
\end{proof}

\begin{proof}[Proof of Theorem~\ref{thm:cycles}]
Proposition~\ref{prop:graphon-bounds} proves
\eqref{eq:cycle-common} and
\eqref{eq:cycle-relative}. If $U\ne0$ on a set of positive
measure, then $t(C_4,U)>0$ by \eqref{eq:c-zero}.
For $k=1$, \eqref{eq:diamond-bound} gives
$16M(K_4-e,W)>2M(P_2,W)\geq1$, making
\eqref{eq:one-apex} strict at $2^{1-2n}$.
For $k=2$, \eqref{eq:fourth-colors} gives $8M(C_4,W)>1$,
so \eqref{eq:two-apices} is strict.
For $k\geq3$, Lemma~\ref{lem:cycle-equality} gives
$M(C_n,W)>2^{1-n}$, making the relative bound strict at
$2^{1-n(k+1)}$.
When $W=1/2$ almost everywhere, every edge factor equals $1/2$
almost everywhere, and equality follows.
\end{proof}

\section{Relation to the Lean formalization}\label{sec:lean}

The Lean libraries \code{TreeApex} and \code{EvenCycleApex}
are developed with Lean~4 \cite{Lean4} and
mathlib \cite{Mathlib}; both dependencies are pinned to version
\code{v4.31.0}. The comparison below uses the theorem statements
and proof dependencies in these libraries.

\subsection{Statements and common foundations}

The principal declarations correspond to the main results as follows.
Names in the first two rows belong to \code{TreeApex}; those in
the remaining rows belong to \code{EvenCycleApex}.
\begin{center}
\small
\begin{tabular}{@{}p{0.34\textwidth}p{0.62\textwidth}@{}}
\toprule
Paper statement & Lean declaration\\
\midrule
Theorem~\ref{thm:trees} &
\code{tree_apex_common}, \code{tree_apex_two_colour}\\[3pt]
Theorem~\ref{thm:tree-one-color} & \code{tree_apex_one_colour}\\[3pt]
Theorem~\ref{thm:cycles}, commonness &
\code{commonality_all_even_all_apices}\\[3pt]
Theorem~\ref{thm:cycles}, relative bound &
\code{apex_relative_of_three_le}, \code{one_le_cycle_normalized}\\[3pt]
Theorem~\ref{thm:cycles}, equality &
\code{commonality_equality_iff_constant}\\
\bottomrule
\end{tabular}
\end{center}
The formal statements allow any probability space $(\Omega,\mu)$,
whereas the paper works on $[0,1]$. Their graphon hypothesis is
joint measurability, pointwise symmetry, and values in $[0,1]$;
the equality conclusion is almost everywhere with respect to
$\mu\otimes\mu$. Lean's \code{pathGraph 3} is our $P_2$, since
its parameter counts vertices. The constants are written as
$4/2^{(k+1)n}$ for trees and $2/2^{n(k+1)}$ for cycles, which
equal the bounds in Theorems~\ref{thm:trees} and~\ref{thm:cycles}.

Both libraries prove finite inequalities for weighted hosts:
nonnegative vertex weights sum to one, and the kernel may have
nonzero diagonal entries. Zero-weight vertices are allowed.
These hosts represent step graphons on finite measurable
partitions. Density continuity and $L^1$ approximation then
transfer closed inequalities to arbitrary graphons, as in the
paper's tree proof. For cycles, the weighted formulation replaces
the uniform dyadic matrices used in
Proposition~\ref{prop:graphon-bounds}. The cycle matrices have
entries $\sqrt{w_iw_j}L(i,j)$, so their traces compute weighted
cycle densities.

\subsection{Organization of the tree proof}

The paper and the formalization use the same relative-entropy
framework on weighted hosts. The definition \code{relEnt} in
\code{TreeApex/Entropy/RelEnt.lean} is
$\mathcal H_\rho(p)$ from \eqref{eq:relative-entropy}, and its
support hypothesis ensures that every logarithm with positive
probability has a positive argument. Both proofs obtain Gibbs'
inequality from $\log u\leq u-1$ and apply it directly to
homomorphism weights.

In \code{Entropy/Triangle.lean}, \code{h_ge} and \code{I_le}
are exactly the two estimates in \eqref{eq:triangle-estimates};
the quantity denoted by \code{A} there is our $\alpha$.
In \code{Entropy/Book.lean}, \code{relEnt_book} establishes
$\mathcal H_b(Q)=\log R+(k-1)h$. The formal page estimate
splits the calculation in \eqref{eq:pages-kl} into two steps:
\code{g_eq} expresses the conditional relative entropy using
the auxiliary ratio
\[
 \Phi(s,x)=\frac{\mu(s,x)R}{\rho_S(s)\nu_s(x)},
\]
and \code{pages_kl} applies Gibbs' inequality to $\mu$ and the
same conditionally independent law $\lambda$ used in the paper.
Together they give the lower bound from
\eqref{eq:pages-kl}; the paper records the divergence explicitly
and combines these steps in a single identity.

The extension in \code{Entropy/TreeLaw.lean} is proved for an
arbitrary symmetric finite law with reference factors on its
vertices and edges, as in Lemma~\ref{lem:tree-extension}. Its
recursion begins with one vertex and adds leaves, whereas the
paper starts from an edge; the first formal edge recovers the
book law. It proves
vertex marginals and support directly, rather than using a
separate theorem asserting every edge marginal. Internally,
vertex $i+1$ has parent $\min(\operatorname{par}(i),i)$.
The theorem \code{exists_recTree_iso} in
\code{Graph/RecTree.lean} identifies every finite tree with
such a recursively labeled tree, so this encoding does not
restrict the tree family. Finally, the two-color construction
is performed on the weighted host $\{0,1\}\times V$, with weights
$w_x/2$ on each copy, before transfer to graphons. This is the
finite version of \eqref{eq:disjoint-colors}. The same generic
extension also proves the tree Sidorenko inequality in
\code{TreeApex/Appendix/Sidorenko.lean}.

\subsection{Cycle certificates and the two-apex case}

The three-apex argument follows the conditional spectral moments,
majority weights, and lifting steps of Section~\ref{sec:cycles}.
Its certificate input consists of exactly three targets:
$Q_3$, $Q_-$, and $Q_+$ from Appendix~\ref{app:certificates}.
In \code{EvenCycleApex/Certificate/Targets.lean}, these are
constructed from their literal edge sets independently of the
certificate matrices. The checker verifies rational
$LDL^{\mathsf T}$ factorizations of the $57$ matrices in these
three certificates. It checks relabeling witnesses only for
masks occurring in the expansions, and accumulates coefficients
in orbit-indexed integers packed in base $2^{128}$.
A proved bound on the total coefficient mass ensures that
equality of packed values implies equality in every orbit;
no completeness claim about a global orbit table is needed.
The concrete checks use \code{decide +kernel}, and their
soundness is proved in Lean. Thus the coefficient identities
and matrix factorizations for these three targets are certified
by proofs evaluated inside the Lean kernel.

The two-apex proof in \code{Finite/TwoApex.lean} omits $Q_2$
entirely. In the paper's density notation, it combines
$\E D_2^2=8M(C_4,W)$,
$\E\Pi_2=128M(C_4^{+1},W)$, the one-apex diamond bound, and
$8M(C_4,W)\leq1+6t(P_2,U)+t(C_4,U)$ to obtain
\begin{equation}\label{eq:lean-two-apex}
 \E V_2^4\geq
 \frac{\bigl(1+2t(P_2,U)+t(C_4,U)/3\bigr)^4}
      {\bigl(1+t(P_2,U)\bigr)^2\bigl(1+6t(P_2,U)+t(C_4,U)\bigr)}
 \geq1.
\end{equation}
The $n/4$ power of this quotient is a lower bound for
$2^{3n-1}M(C_n^{+2},W)$ for every even $n\geq4$.
The scalar inequality is verified by expanding the difference
into nonnegative terms. Moreover, the quotient exceeds one when
$t(C_4,U)>0$,
which suffices for rigidity. This is a substantive change in an
auxiliary conclusion: the Lean development does not establish
$\E\Pi_2\geq\E D_2^2$ or the paper's stronger two-apex bound
\eqref{eq:two-apices}.
It establishes the same commonness and equality conclusions
with two apices. In particular, the formal theorem
\code{three_universal_graph_inequalities} covers three of the
four inequalities in Lemma~\ref{lem:cert-input}.

\subsection{Rigidity and verification scope}

The cycle rigidity proof also follows a different analytic route.
Instead of applying the compact self-adjoint spectral theorem
directly as in \eqref{eq:infinite-trace}, Lean first proves a
finite spectral remainder estimate. For even $n>4$, it transfers
the continuous inequality
\begin{align*}
 8M(C_4,W)
 &\leq\frac12\sum_{\varepsilon=\pm1}
 \Bigl[t(C_n,S_\varepsilon)^{4/n}\\
 &\qquad+4^{(n-4)/(n-2)}
   \max\{t(C_n,S_\varepsilon)
        -(1+\varepsilon t(K_2,U))^n,0\}^{2/(n-2)}\Bigr]
\end{align*}
from weighted hosts to graphons. When $M(C_n,W)=2^{1-n}$,
the mean bound forces $t(K_2,U)=0$ and both signed color cycle
densities equal one, so this estimate gives $8M(C_4,W)=1$
and hence $t(C_4,U)=0$. In \code{Equality/CZero.lean},
this vanishing first gives $U\circ U=0$ almost everywhere. Symmetry
then makes every bilinear form of $U$ against bounded product
functions vanish; approximation by finite sums of such products
in $L^2$ yields $U=0$ almost everywhere. This replaces the
infinite-dimensional spectral argument while preserving its
conclusion.

The main graphon inequalities and the cycle equality
characterization are thus formalized, with the two-apex auxiliary
bound replaced as explained above. The tree assertions of equality
at constant graphons follow from \code{homDensity_const}, but
are not stated as separate tree headline theorems.
The random-sampling lemma and the equivalence with finite
Ramsey multiplicity in Appendix~\ref{app:finite} are not part
of either Lean development. The audit files
\code{CheckAxioms.lean} and \code{CheckAxiomsTree.lean} check
the dependencies of the main results and their supporting
lemmas; their outputs contain only \code{propext},
\code{Classical.choice}, and \code{Quot.sound}, or subsets
thereof. In particular, the audited results depend on neither
unproved placeholders nor an axiom for native computation.
The checks \code{CheckStatementsTree.lean} and
\code{CheckImportsTree.lean} additionally verify the tree
statement signatures and the absence of cycle certificate
modules from the tree library's import closure.

\section{Further directions}

Trees and even cycles have the Sidorenko property, so the two
theorems settle Conjecture~\ref{conj:apices} for these base
families. Their proofs isolate different ways in which the base
structure can be exploited. On a tree, consistency of local edge
laws suffices to construct a global homomorphism distribution.
On a cycle, the closing edge is instead handled by an even
spectral trace. A proof for more general connected Sidorenko
bases would need to control the compatibility of the additional
edges with the chosen local distributions or moments.

The relative comparison for cycles is proved for $k\geq3$.
The bound \eqref{eq:two-apices} suffices for commonness and
rigidity with two apices, but does not imply
$M(C_n^{+2},W)\geq2^{-2n}M(C_n,W)$ for every graphon.
The sharp tree inequality and the weighted
fourth-moment comparison are stronger than their commonness
consequences and may be useful for related density problems.

\appendix
\section{Exact flag algebra certificates}\label{app:certificates}

We specify the four certificate identities used in
Lemma~\ref{lem:cert-input}. The presentation uses signed
homomorphism densities rather than induced colored densities;
the positivity mechanism is the rooted quadratic-form principle
of flag algebras \cite{Razborov}. All four certificates use the
common feature schema described below.

\subsection{Targets in the signed graph algebra}

Label six formal vertices $0,\ldots,5$ and order their fifteen
unordered pairs lexicographically. A graph is encoded by the
integer mask of its edge set; $[E]$ denotes its formal basis
element. Evaluating a linear combination means replacing
$[E]$ by the homomorphism density of that graph in $U$.
Unused vertices are isolated and cause no change in evaluation.
For an edge set $E$, put
\[
 \mathsf P_\eta(E)=
 \sum_{\substack{A\subseteq E\\|A|\equiv\eta\pmod2}}[A]
 \qquad(\eta=0,1).
\]
If two basis graphs have disjoint edge sets, their labeled product
is defined by taking the union. We only use this product when no
edge repeats.

For $s=2,3$, let $I_s=\{s,s+1,s+2\}$ and define
\begin{align*}
 E_s&=\{ij:0\leq i<s,\ j\in I_s\}
                    \cup\{s(s+1),s(s+2)\},\\
 B_s&=\{ij:0\leq i<s,\ j\in\{s+1,s+2\}\}.
\end{align*}
Here $ij$ denotes the unordered pair $\{i,j\}$.
The two mean targets are
\[
 Q_s=\mathsf P_0(E_s)-\mathsf P_0(B_s),\qquad s=2,3.
\]
Their evaluations are $\E\Pi_s-\E D_s^2$, by expanding the conditional
integrals. For the weighted targets, set
\[
F_\eta=2\mathsf P_\eta(E_3)-\mathsf P_\eta(B_3),\quad
 B=[01]+[02]+[12]-[01,02,12].
\]
Define the two weighted targets directly in this graph basis:
\begin{align}
 Q_-={}&3F_0-3\mathsf P_0(\{01,12,23,03\})+F_1-BF_0\notag\\
 &\quad-13[01]-[01,02,12]-12[01,23]\notag\\
 &\quad+4[01,23,24,34]-12[01,12,23],
 \label{eq:target-minus}\\
 Q_+={}&6F_0-6\mathsf P_0(\{01,12,23,03\})+F_1+BF_0\notag\\
 &\quad-19[01]-23[01,02,12]+12[01,23]\notag\\
 &\quad-4[01,23,24,34]-24[01,12,23].
 \label{eq:target-plus}
\end{align}
The product $BF_0$ is legitimate: $B$ uses only pairs among
$0,1,2$, while $E_3$ and $B_3$ have no such pairs.
Evaluation sends $F_0,F_1,BF_0$ to
$\E F,\E\varepsilon F,\E\maj F$, respectively, and sends
$\mathsf P_0(\{01,12,23,03\})$ to $8M(C_4,W)$.
The basis graph $[01,23,24,34]$ is a disjoint edge and triangle,
whose density is $t(K_2,U)t(K_3,U)$.
Consequently \eqref{eq:target-minus} and
\eqref{eq:target-plus} evaluate to the left sides minus the
right sides of \eqref{eq:G} and \eqref{eq:H}.

\subsection{Rooted quadratic forms}

Choose $r$ roots and $\ell$ further vertices. A feature graph $A$
has no root--root edges. Its rooted density is
\[
 \phi_A(\mathbf z)=\int
           \prod_{ij\in E(A)}U(x_i,x_j)\dd\mathbf x,
\]
where the root coordinates are fixed at $\mathbf z$ and the other
coordinates are integrated independently. For a root type
$t\subseteq\binom{\{0,\ldots,r-1\}}2$, let
\[
 p_t(\mathbf z)=2^{-\binom r2}
       \prod_{e\in t}(1+U_e)\prod_{e\notin t}(1-U_e).
\]
Every factor is nonnegative when $|U|\leq1$.
For a feature vector $\phi$ and a positive semidefinite matrix $Q$,
the expression
\begin{equation}\label{eq:rooted-sos}
 \int p_t(\mathbf z)\phi(\mathbf z)^{\mathsf T}
                      Q\phi(\mathbf z)\dd\mathbf z
\end{equation}
is therefore nonnegative. This remains true for $r=0$, when the
root space has one point and $p_t=1$.

Each certificate uses nineteen feature blocks, ordered as follows.
The root-type integers are binary masks in the lexicographic
root-pair order.
\begin{center}
\begin{tabular}{@{}cccl@{}}
\toprule
Roots $r$ & Free vertices $\ell$ & Root types $t$ & Block dimensions\\
\midrule
0 & 3 & 0 & 3\\
1 & 2 & 0 & 5\\
2 & 2 & 0, 1 & 19, 19\\
3 & 1 & 0, 1, 3, 7 & four blocks of 7\\
4 & 1 & $0,1,3,7,11,12,13,15,30,31,63$ & eleven blocks of 15\\
\bottomrule
\end{tabular}
\end{center}
For $r=0$, the features have edge sets $\{01\}$,
$\{01,02\}$, and $\{01,02,12\}$, in that order.
For $r=1,2$, order the permitted edges as
\[
 (0,r),(0,r+1),\ldots,(r-1,r),(r-1,r+1),(r,r+1).
\]
Enumerate nonempty subsets by increasing binary mask and retain
a mask exactly when it is at most its image under swapping the
two free vertices. These are the five and nineteen features
specified in the table. For $r=3,4$, use every nonempty subset
of the edges $(i,r)$, $0\leq i<r$, in increasing mask order.

The common matrix denominator is
\[
 \mathscr D=90315258984881964711936.
\]
For each target, the certificate consists of nineteen integer
matrices $A_j$ in this block order; the quadratic-form matrices
are $A_j/\mathscr D$.

\subsection{The finite identities and their verification}

For a block with roots $r$ and $\ell$ free vertices, let
$A_a'$ be a copy of feature graph $A_a$ using a disjoint set of
free vertices and the same roots. Expanding \eqref{eq:rooted-sos}
shows that its formal polynomial, after multiplication by
$64\mathscr D$, is
\begin{equation}\label{eq:block-expansion}
 \mathcal R_j:=2^{6-\binom r2}\sum_{a,b}(A_j)_{ab}
 \sum_{T\subseteq\binom{\{0,\ldots,r-1\}}2}
    (-1)^{|T\setminus t_j|}
       [T\cup E(A_a)\cup E(A_b')].
\end{equation}
The root edges and the two free extensions are edge-disjoint,
which makes every monomial a simple graph. Independent integration
of the two extensions is exactly multiplication of their rooted
features.

Let $\mathsf N$ relabel each six-vertex mask to an isomorphism
representative and extend this map linearly. Since densities are
invariant under relabeling, equality after applying $\mathsf N$
implies equality of evaluations for every symmetric kernel.
The exact certificate identity for each
$Q\in\{Q_2,Q_3,Q_-,Q_+\}$ is
\begin{equation}\label{eq:certificate-identity}
 \mathsf N(64\mathscr D Q)=
 \mathsf N\left(\sum_{j=1}^{19}\mathcal R_j\right).
\end{equation}
Moreover, every integer matrix $A_j$ has a rational factorization
$A_j=L_j\diag(\delta_j)L_j^{\mathsf T}$ with $L_j$ unit lower
triangular and all entries of $\delta_j$ strictly positive.
Consequently $A_j/\mathscr D$ is positive definite, since
\[
 x^{\mathsf T}(A_j/\mathscr D)x
   =\mathscr D^{-1}\sum_i(\delta_j)_i
                       ((L_j^{\mathsf T}x)_i)^2\geq0.
\]
Combining these factorizations, \eqref{eq:block-expansion}, and
\eqref{eq:certificate-identity} expresses each target evaluation
as a sum of nineteen integrals of the form
\eqref{eq:rooted-sos}. This proves Lemma~\ref{lem:cert-input}.

The finite arithmetic in \eqref{eq:certificate-identity} has
been verified by two independent exact-arithmetic calculations.
One expands the targets and rooted quadratic forms, compares
their coefficients on six-vertex isomorphism classes, and checks
positive definiteness by exact leading principal minors.
The other constructs explicit relabeling witnesses using
adjacent transpositions and expands the polynomials by a separate
convolution procedure. It checks all $32768$ relabeling witnesses
and all $131072$ coefficient positions across the four targets,
including zeros. For the $76$ matrices, rational
$LDL^{\mathsf T}$ factorizations check $956$ positive pivots
and reconstruct all $13708$ matrix entries exactly.
The weighted-comparison identities and normalization calculations
are also verified by exact arithmetic. The Lean verification
of the three targets $Q_3,Q_-,Q_+$ is described in
Section~\ref{sec:lean}; the two-apex Lean argument uses the
alternative bound \eqref{eq:lean-two-apex}.

\section{The finite coloring formulation}\label{app:finite}

We record why \eqref{eq:common} gives the usual Ramsey
multiplicity statement. Fix a simple graph $F$ with $v$ vertices
and $e\geq1$ edges. From a red--blue coloring of $K_N$, form its
red adjacency step graphon on $N$ equal intervals, assigning
arbitrary colors on the $N$ diagonal cells. Then $N^vM(F,W)$
counts monochromatic homomorphisms into the resulting looped
host. At most $\binom v2N^{v-1}$ vertex maps have a collision.
Each map can be monochromatic in at most one color, since $F$
has an edge. Thus commonness implies at least
\[
 2^{1-e}N^v-\binom v2N^{v-1}
\]
injective labeled monochromatic copies in every coloring.
Conversely, the expected number in a fair independent coloring is
$(N)_v2^{1-e}$, so some coloring has at most this many. Dividing
by $N^v$ and letting $N\to\infty$ identifies the asymptotic
minimum as $2^{1-e}$. The edgeless case is immediate separately.

Conversely, sample a host from a graphon $W$ as in
Lemma~\ref{lem:finite-approx} and color its present edges red and
its absent edges blue. The same expectation and variance calculation
applies simultaneously to the two colors, so their normalized
injective counts converge to $t(F,W)$ and $t(F,1-W)$.
An asymptotic lower bound of $2^{1-e}$ for every finite coloring
therefore forces $M(F,W)\geq2^{1-e}$. This proves the equivalence
used in the introduction.

\bibliographystyle{plainurl}
\bibliography{apices_commonness}
\end{document}
